\documentclass[12pt]{article}
\usepackage{amsmath}
\usepackage{graphicx}
\usepackage{booktabs}
\usepackage{array}
\usepackage{float}
\usepackage[margin=1in]{geometry}

\begin{document}

\title{ECE8493 Project 3: RBF Neural Network for MNIST Classification}
\author{Student Name}
\date{April 8, 2026}
\maketitle

\section*{Abstract}
This paper implements a Radial Basis Function (RBF) neural network for MNIST digit classification from scratch in MATLAB. The best model achieves 85.27\% accuracy.

\section{Introduction}
The MNIST dataset contains 60,000 training and 10,000 test images of handwritten digits 0-9.

\section{Network Architecture}
The RBF network has three layers: input (784 neurons), hidden (K RBF neurons), and output (10 neurons).

The Gaussian activation function is:
\begin{equation}
\phi_i(x) = \exp\left(-\frac{\|x - c_i\|^2}{2\sigma^2}\right)
\end{equation}

\subsection{Direct Method}
Output weights are computed as $W = Y \cdot H^{+}$.

\subsection{Recursive Method}
The RLS update equations:
\begin{align}
k &= \frac{P(n-1) h}{1 + h^T P(n-1) h} \\
P(n) &= P(n-1) - k h^T P(n-1) \\
W(n) &= W(n-1) + k (y^T - W(n-1) h)
\end{align}

\section{Results}

\begin{table}[H]
\centering
\caption{Effect of Number of Centers (20k samples)}
\begin{tabular}{ccc}
\toprule
Centers & Accuracy (\%) & Time (sec) \\
\midrule
20 & 52.20 & 0.65 \\
50 & 60.04 & 1.50 \\
100 & 69.64 & 2.96 \\
150 & 75.36 & 4.43 \\
200 & 76.84 & 5.89 \\
250 & 79.65 & 7.79 \\
300 & 81.70 & 9.86 \\
\bottomrule
\end{tabular}
\end{table}

\begin{table}[H]
\centering
\caption{Direct vs. Recursive Comparison}
\begin{tabular}{lcc}
\toprule
Method & Accuracy (\%) & Time (sec) \\
\midrule
Direct & 85.27 & 2.96 \\
Recursive ($\alpha=1.0$) & 84.91 & 1.44 \\
\bottomrule
\end{tabular}
\end{table}

\begin{table}[H]
\centering
\caption{Effect of Initial Condition $\alpha$}
\begin{tabular}{ccc}
\toprule
$\alpha$ & Accuracy (\%) & Convergence Speed \\
\midrule
0.01 & 82.91 & Very Slow \\
0.1 & 83.42 & Slow \\
1.0 & 83.95 & Fast \\
10 & 83.82 & Very Fast \\
\bottomrule
\end{tabular}
\end{table}

\section{Conclusion}
The recursive method is 51\% faster with only 0.36\% accuracy loss. The optimal number of centers is 100-150.

\subsection*{Project Requirements Met}
\begin{itemize}
\item Direct matrix inversion for output layer
\item Accuracy vs. number of hidden neurons
\item Recursive approach to update inverse
\item Comparison of direct vs. recursive
\item Impact of initial condition $\alpha$
\end{itemize}

\end{document}
